% ============================================================================
\section{Experimental Setup}
\label{sec:experiments}
% ============================================================================

We distill four specialist teachers, for mathematics, code, instruction
following, and science, into one Qwen3-1.7B student. Table~\ref{tab:pipeline}
in the appendix lists the models, data, hyperparameters, and system
configuration; entries in red remain to be filled from the final runs.

\subsection{Data and budget}

Each task has a fixed stream of $16{,}000$ prompts, shuffled once and consumed
without replacement. All methods read a task's stream in the same order, so
the $n$th prompt a method requests from a task is the same prompt for every
method; the allocation determines only how far into each stream a method has
progressed at a given step. Training runs for $500$ steps of $64$ responses,
$32{,}000$ responses in total. A task at the maximum count on every step
consumes exactly $500\times8\times4=16{,}000$ prompts, so no prompt is ever
reused. The objective weights are the normalized inverse initial task losses,
measured once at the initial student and shared by all runs.

\subsection{Methods}

Uniform assigns $(4,4,4,4)$ micro-batches on every step. GPAS uses
Equation~\ref{eq:batch_gpas_allocation} and Cost-GPAS minimizes
Equation~\ref{eq:batch_cost_objective}. RawNoise is a diagnostic that applies
the GPAS rule to the unscaled noise $\mathbb E\|g_i-\mathbb E[g_i]\|_2^2$
instead of $e_i$, isolating the effect of the optimizer coordinates. All
methods share the optimizer, learning-rate schedule, count bounds, and prompt
streams, and every method starts with one Uniform step. Uniform and GPAS use
three seeds each; the seed changes the prompt shuffle and the response
sampling, and streams stay matched across methods within a seed. Cost-GPAS
and RawNoise use one seed. \missing{Run RawNoise only if the held-out
variance check shows that raw and scaled allocations differ materially.}

\subsection{Measurements}

\paragraph{Held-out teacher loss.} The primary training metric is the
weighted teacher loss $\sum_i w_iL_i$ on a fixed held-out set of $128$ prompts
per task, evaluated from checkpoints every $50$ steps with a fixed sampling
seed. Training-batch losses are reported in the appendix because they depend
on which prompts each method drew. The common loss target for compute
comparisons is Uniform's final held-out loss.

\paragraph{Compute.} GPU hours are full wall time times the number of reserved
GPUs, including teacher loading and waiting. Each step's trace records
per-task rollout, scoring, backward, and teacher-loading time, from which
$\tau_i$ and $C$ are obtained (Appendix~\ref{app:system_details}).

\paragraph{Held-out gradient variance.} At early, middle, and late Uniform
checkpoints, we generate $32$ fresh micro-batches per task, estimate
$\widehat e_i$ on each half, and evaluate the allocation each rule computes
from one half on the noise measured on the other half
(Appendix~\ref{app:heldout}). This tests whether the within-step estimate
predicts fresh noise and how much variance each rule removes, independently
of the online ratio $H$.

\paragraph{Capability.} We report MATH-500 greedy pass@1, a fixed
LiveCodeBench pass@1 slice, IFBench strict accuracy, and GPQA-Diamond
average@4 for the initial student, each teacher, and each method, together
with the per-task normalized gain
$(s-s_{\rm init})/(s_{\rm teacher}-s_{\rm init})$ and its average.
\missing{Confirm benchmark versions and decoding settings; add AIME24/25
average@16 or a second science benchmark if the initial student leaves too
little headroom.} We also report mean response length and truncation rate per
task.

\subsection{System}

The target hardware holds one teacher at a time, so each step loads the four
teachers in a fixed cyclic order and every method pays the same loading cost.
Generation stops at $8{,}192$ tokens per response and truncated responses are
kept, because the token-level loss needs no completed response. The cap
bounds the generation time of each task block, so the step time is
predictable from the counts through $C+\sum_i m_i\tau_i$. We report the
overhead ratio $C/\sum_i m_i\tau_i$ next to GPU hours, because a large ratio
limits what Cost-GPAS can save.

% ============================================================================
\section{Results}
\label{sec:results}
% ============================================================================

\subsection{Noise and allocation dynamics}

\begin{figure}[H]
    \centering
    \missingfigure[0.18\textheight]{Insert Figure 2 from the four-teacher
    logs: (a) scaled and raw noise by task over training; (b) GPAS and
    Cost-GPAS counts; (c) variance ratio $H$ with the bound $2$; (d) per-task
    time $\tau_i$ and the overhead ratio.}
    \caption{Optimizer-scaled noise and the resulting allocations.
    \missing{State the main quantitative observation in the caption.}}
    \label{fig:allocation_dynamics}
\end{figure}

\begin{table}[H]
\centering
\small
\caption{Held-out gradient variance relative to Uniform at three Uniform
checkpoints, for the allocation each rule computes from the other half of the
fresh micro-batches. \missing{Fill from the checkpoint statistics; add the
relative error of $\widehat e_i$ against the held-out half.}}
\label{tab:heldout_variance}
\begin{tabular}{lccc}
\toprule
Allocation & Early & Middle & Late \\
\midrule
Uniform & $1.00$ & $1.00$ & $1.00$ \\
RawNoise & & & \\
GPAS & & & \\
Cost-GPAS & & & \\
\bottomrule
\end{tabular}
\end{table}

\missing{Report the cross-task range and drift of $\widehat e_i$, the fraction
of steps in which raw and scaled noise rank the tasks differently, how often
each task sits at a count bound, and the median and quartiles of $H$.}

\subsection{Training efficiency}

\begin{figure}[H]
    \centering
    \missingfigure[0.17\textheight]{Insert Figure 3: (a) held-out weighted
    teacher loss versus optimizer step for Uniform and GPAS, with seed bands;
    (b) held-out weighted teacher loss versus GPU hours for Uniform, GPAS, and
    Cost-GPAS. Per-task curves go to the appendix.}
    \caption{Held-out teacher loss at matched steps and at matched compute.
    \missing{State the matched-step loss difference and the GPU-hour saving at
    Uniform's final loss.}}
    \label{fig:training_efficiency}
\end{figure}

\subsection{Capability and compute}

\begin{table}[H]
\centering
\scriptsize
\setlength{\tabcolsep}{2.5pt}
\caption{Four-teacher capability and efficiency. \missing{Populate from the
evaluation outputs and timing logs; state every task-level difference from
Uniform next to the average.}}
\label{tab:main_results}
\begin{tabular}{lccccccc}
\toprule
& Math & Code & IF & Science & Norm. avg. & Final held-out $F$ & GPU h to target \\
\midrule
Initial student & & & & & & & \\
Teachers & & & & & & & \\
Uniform & & & & & & & \\
GPAS & & & & & & & \\
Cost-GPAS & & & & & & & \\
RawNoise & & & & & & & \\
\bottomrule
\end{tabular}
\end{table}
